\section{Proof of the Fisher information theorem}\label{app:fisher}

We keep the notation of Section~\ref{sec:fisher}. So $0<q<1$, $p=1-q$,
$\rho=1-2q$, $u=q/p$, $x=Nq$ and $y=Nu=x/p$. Let $\sigma_i\in\{-1,1\}$ be the
direction of step $i$, with $\sigma_i=1$ for a step to the right. Then
$X=\sigma_1+\dots+\sigma_N$, and $J$ is the number of $i<N$ with
$\sigma_{i+1}\ne\sigma_i$. Put $Z=\E[J\mid K]$. By
Proposition~\ref{prop:eta}, $e_N=\Var Z/\Var J$.

For an interior bin, Paper~A's Bessel estimate makes $Z$ close to
$y\sqrt{1-W^2}+\frac12$. When $W$ is small this is close to the quadratic
$y(1-W^2/2)+\frac12$, and $W$ is of order $x^{-1/2}$. So $Z$ is close to its
projection on the span of $1$ and $X^2$, and $e_N$ is close to $R^2_N$. We
first prove the exact moments behind the closed form of $R^2_N$.

\paragraph{Exact moments.}

\begin{lemma}[exact moments]\label{lem:fisher-moments}
For $N\ge1$ and $0<q<1$,
\begin{align*}
 \E X^2&=\frac{Np}q-\frac{\rho(1-\rho^N)}{2q^2},\\
 \Cov(J,X^2)&=-\frac p{q^2}\bigl[Nq(1+\rho^N)-p(1-\rho^N)\bigr],\\
 \E X^4&=\frac{3N^2p^2}{q^2}-\frac{Np(1+10\rho+\rho^2)}{2q^3}
 +\frac{\rho(1+2\rho)(2+\rho)(1-\rho^N)}{2q^4}
 -\frac{3Np\rho^{N+1}}{q^3},
\end{align*}
and $\Var X^2=2N^2p^2V/q^2$, with $V$ as in Theorem~\ref{thm:fisher}.
\end{lemma}

\begin{proof}
\emph{Step 1.} For real $\phi$ and $\chi$ put
$m_N(\phi,\chi)=\E e^{\phi X+\chi J}$. Let $A$ be the $2\times2$ matrix with
entries $A_{ab}=\Prob(\sigma_2=b\mid\sigma_1=a)\,e^{\phi b+\chi\mathbf
1[a\ne b]}$ for $a,b\in\{-1,1\}$, and let $v_a=\frac12e^{\phi a}$. Then
$m_N=v^\top A^{N-1}\ones$ with $\ones=(1,1)^\top$. So, as formal power
series in $z$,
\[
 \sum_{N\ge1}m_Nz^N=zv^\top(I-zA)^{-1}\ones=\frac{U(z)}{D(z)},\qquad
 U(z)=zv^\top\operatorname{adj}(I-zA)\ones,\quad D(z)=\det(I-zA).
\]
Both $U$ and $D$ are polynomials in $z$ of degree at most $2$, and their
coefficients are smooth functions of $(\phi,\chi)$. Let $\partial$ be a
partial derivative in $(\phi,\chi)$ of total order $m$. By induction on $m$
and the quotient rule, $\partial(U/D)=U_\partial/D^{m+1}$, where $U_\partial$
is a polynomial in $z$ of degree at most $2m+2$. Indeed each derivative
multiplies the numerator by $D$ or by a derivative of $D$, and this raises
its degree by at most $2$. At $\phi=\chi=0$ the matrix $A$ has
trace $2p=1+\rho$ and determinant $p^2-q^2=\rho$, so
$D(z)=(1-z)(1-\rho z)$. Now let $q\ne\frac12$. Then $\rho\notin\{0,1\}$, and
the roots $1$ and $1/\rho$ of $D$ are distinct. Partial fractions write
$U_\partial/D^{m+1}$ as a constant plus a combination of $(1-z)^{-j}$ and
$(1-\rho z)^{-j}$ with $1\le j\le m+1$. The constant only changes the
coefficient of $z^0$. The coefficient of $z^N$ in $(1-z)^{-j}$ is
$\binom{N+j-1}{j-1}$, a polynomial in $N$ of degree $j-1$, and in
$(1-\rho z)^{-j}$ it is the same polynomial times $\rho^N$. Hence, as a
sequence in $N\ge1$, the derivative $\partial m_N$ at $\phi=\chi=0$ lies in
\[
 \mathcal S_m=\operatorname{span}\{N^i,\ N^i\rho^N:0\le i\le m\},
\]
which has dimension at most $2m+2$.

\emph{Step 2.} Every sequence in $\mathcal S_m$ satisfies the linear
recurrence with characteristic polynomial $(t-1)^{m+1}(t-\rho)^{m+1}$, which
has order $2m+2$. So a sequence in $\mathcal S_m$ is determined, for
$N\ge1$, by its values at $N=1,\dots,2m+2$. In particular, if these values
are $0$, the sequence is $0$ for every $N\ge1$.

\emph{Step 3.} The moments $\E X^2$, $\E[JX^2]$ and $\E X^4$ are the
derivatives $\partial_\phi^2m_N$, $\partial_\chi\partial_\phi^2m_N$ and
$\partial_\phi^4m_N$ at $0$, so they lie in $\mathcal S_4$. The claimed
formulas for $\E X^2$ and $\E X^4$ lie in $\mathcal S_2\subset\mathcal S_4$.
Since $\E J=(N-1)q$, the claimed formula for
$\E[JX^2]=\Cov(J,X^2)+(N-1)q\,\E X^2$ also lies in $\mathcal S_2$. Hence the
difference between each moment and its claimed formula lies in
$\mathcal S_4$, which has dimension at most $10$. By Step~2 it is enough to
check that the difference vanishes at $N=1,\dots,10$. This is a finite exact
computation. The script \texttt{verify\_sticky\_fisher.py} computes the 3
moments exactly for $N=1,\dots,12$ by a dynamic program over paths, with $q$
kept as a symbol, and every difference is identically $0$. This proves the 3
formulas for $q\ne\frac12$. For fixed $N$ both sides are rational functions
of $q$ with no pole at $q=\frac12$, so the case $q=\frac12$ follows by
continuity.

\emph{Step 4.} We have $\Var X^2=\E X^4-(\E X^2)^2$. With the formulas of
Step~3, the claim $\Var X^2=2N^2p^2V/q^2$ becomes an identity between
rational functions of 3 independent variables $N$, $q$ and $R$, where $R$
stands for $\rho^N$ and $\rho^{N+1}=\rho R$. The same script checks this
identity with sympy.
\end{proof}

\begin{proof}[Proof of the closed form in Theorem~\ref{thm:fisher}]
Let $N\ge2$. Then $X^2$ takes the value $N^2$ on the constant paths and the
smaller value $(N-2)^2$ on the paths that switch only at the last step. So
$\Var X^2>0$, and $V>0$ by Lemma~\ref{lem:fisher-moments}. The same lemma
gives
\[
 \Cov(J,X^2)=-\frac{Np}q\Bigl[1+\rho^N-\frac{p(1-\rho^N)}x\Bigr],
\]
so $\Cov(J,X^2)^2/\Var X^2$ is the squared bracket divided by $2V$. Dividing
by $\Var J=(N-1)pq$ gives the formula for $R^2_N$.
\end{proof}

\paragraph{Three lemmas.}

\begin{lemma}[convexity]\label{lem:fisher-convex}
Let $\delta>0$, and let $f$ and $g$ be convex and differentiable on
$[t-\delta,t+\delta]$ with $0\le g-f\le\eta$ there. Suppose also that $g$ is
twice differentiable with $g''\le\kappa$ on this interval. Then
\[
 \lvert f'(t)-g'(t)\rvert\le\frac\eta\delta+\frac{\kappa\delta}2 .
\]
\end{lemma}

\begin{proof}
By the convexity of $f$ and then the bounds on $g-f$, we have
$f'(t)\le(f(t+\delta)-f(t))/\delta\le(g(t+\delta)-g(t)+\eta)/\delta$. By
Taylor's theorem $g(t+\delta)\le g(t)+\delta g'(t)+\kappa\delta^2/2$, so
$f'(t)\le g'(t)+\eta/\delta+\kappa\delta/2$. In the same way
\[
 f'(t)\ge\frac{f(t)-f(t-\delta)}\delta\ge\frac{g(t)-g(t-\delta)-\eta}\delta
 \ge g'(t)-\frac\eta\delta-\frac{\kappa\delta}2 .\qedhere
\]
\end{proof}

\begin{lemma}[a Bessel ratio]\label{lem:fisher-bessel}
For $t>0$ let $R(t)=I_1(t)/I_0(t)$. Then
\[
 \frac{\sqrt{1+t^2}-1}t\le R(t)<1 .
\]
In particular $0<1-R(t)\le1/t$.
\end{lemma}

\begin{proof}
Since $I_0'=I_1$ and $I_1'(t)=I_0(t)-I_1(t)/t$, the function $R$ satisfies
$R'=1-R/t-R^2$. From the power series, $R>0$ and
$R(t)=t/2-t^3/16+O(t^5)$ as $t\to0$. If $R(t_0)=1$ for a first $t_0>0$, then
$R'(t_0)\ge0$, but the equation gives $R'(t_0)=-1/t_0$. So $R<1$.

Put $s=\sqrt{1+t^2}$ and $\ell(t)=(s-1)/t=t/(s+1)$. A short computation
gives $1-\ell/t-\ell^2=\ell/t=1/(s+1)$ and $\ell'=1/(s(s+1))$. Hence
$\ell'-(1-\ell/t-\ell^2)=(1-s)/(s(s+1))\le0$. For $d=R-\ell$ this gives
\[
 d'\ge-\frac dt-(R^2-\ell^2)=-d\Bigl(\frac1t+R+\ell\Bigr).
\]
Let $\Theta(t)=t\exp\int_0^t(R+\ell)$. Then
$(\Theta d)'=\Theta\bigl(d'+d(1/t+R+\ell)\bigr)\ge0$. Since
$\ell(t)=t/2-t^3/8+O(t^5)$, we have $d(t)=O(t^3)$, so $\Theta(t)d(t)\to0$ as
$t\to0$. Hence $\Theta d\ge0$ and $d\ge0$. Finally
$1-\ell(t)=(t+1-s)/t\le1/t$, since $s\ge t$.
\end{proof}

\begin{lemma}[moments of the endpoint]\label{lem:fisher-W}
\begin{enumerate}[(a)]
\item For each $m\ge1$ there is a constant $C_m$ such that
$\E W^{2m}\le C_mx^{-m}$ whenever $0<q<1$ and $x\ge1$.
\item If $0<q\le\frac12$, then $\E W^4\le3p^2/x^2+9/(2x^4)$.
\end{enumerate}
\end{lemma}

The reason for the order $x^{-m}$ is that $X$ is roughly a sum of about $x$
runs with alternating signs and lengths of order $1/q$. So $X$ is of order
$\sqrt x/q$ and $W$ is of order $x^{-1/2}$.

\begin{proof}
(a) Extend the walk by a further step $\sigma_{N+1}$ with the same rule, and
let $\mathcal F_i$ be generated by $\sigma_1,\dots,\sigma_i$. Since
$\E[\sigma_{i+1}\mid\mathcal F_i]=(p-q)\sigma_i=\rho\sigma_i$, the variables
$\xi_{i+1}=\sigma_{i+1}-\rho\sigma_i$ are martingale differences. Summing
$\sigma_{i+1}-\sigma_i=-2q\sigma_i+\xi_{i+1}$ over $i=1,\dots,N$ gives
\[
 X=\frac{\sigma_1-\sigma_{N+1}+\Xi}{2q},\qquad
 \Xi=\sum_{i=2}^{N+1}\xi_i .
\]
If step $i+1$ repeats step $i$, then $\xi_{i+1}^2=(1-\rho)^2=4q^2$.
Otherwise $\xi_{i+1}^2=(1+\rho)^2=4p^2\le4$. So
\[
 \sum_{i=2}^{N+1}\xi_i^2\le4q^2N+4J',
\]
where $J'\sim\operatorname{Bin}(N,q)$ is the number of switches among the
$N+1$ steps. Burkholder's inequality~\cite[Ch.~2]{hallheyde1980} bounds the
moment of order $2m$ of a martingale by a constant $c_m$ times the $m$-th
moment of the sum of its squared differences. So
\[
 \E\,\Xi^{2m}\le c_m\E\bigl(4q^2N+4J'\bigr)^m
 \le c_m8^m\bigl((q^2N)^m+\E J'^m\bigr).
\]
Here $q^2N=qx\le x$. The factorial moments of $J'$ are
$\E[J'(J'-1)\cdots(J'-j+1)]=N(N-1)\cdots(N-j+1)q^j\le x^j$, and $J'^m$ is a
combination of the falling factorials $J'(J'-1)\cdots(J'-j+1)$, $1\le j\le
m$, with nonnegative coefficients (the Stirling numbers of the second kind).
So $\E J'^m\le b_mx^m$ for $x\ge1$, where $b_m$ is the sum of these
coefficients. Hence $\E\,\Xi^{2m}\le c_m8^m(1+b_m)x^m$. Since
$\lvert\sigma_1-\sigma_{N+1}\rvert\le2$ and $x\ge1$,
\[
 \E X^{2m}\le(2q)^{-2m}2^{2m-1}\bigl(2^{2m}+\E\,\Xi^{2m}\bigr)
 \le C_mq^{-2m}x^m ,
\]
and $\E W^{2m}=N^{-2m}\E X^{2m}\le C_mx^{-m}$.

(b) Dividing the formula for $\E X^4$ in Lemma~\ref{lem:fisher-moments} by
$N^4$ gives
\[
 \E W^4=\frac{3p^2}{x^2}-\frac{p(1+10\rho+\rho^2)}{2x^3}
 +\frac{\rho(1+2\rho)(2+\rho)(1-\rho^N)}{2x^4}-\frac{3p\rho^{N+1}}{x^3}.
\]
For $q\le\frac12$ we have $0\le\rho<1$. Then the second and the last terms
are at most $0$, and the third is at most $1\cdot3\cdot3/(2x^4)$.
\end{proof}

\paragraph{Proof of the bound in Theorem~\ref{thm:fisher}.}
Let $\mathcal P$ be the orthogonal projection in $L^2$ onto the span of $1$
and $X^2$. Since $X^2$ is a function of $K$, we have $\E Z=\E J$ and
$\Cov(Z,X^2)=\Cov(J,X^2)$. So
$\mathcal PZ=\E Z+\Cov(J,X^2)(X^2-\E X^2)/\Var X^2$, and $Z-\mathcal PZ$ is orthogonal
to $1$ and to $X^2$. Hence
\[
 \Var Z=\frac{\Cov(J,X^2)^2}{\Var X^2}+\E(Z-\mathcal PZ)^2 ,
\]
and dividing by $\Var J$ gives $e_N-R^2_N=\E(Z-\mathcal PZ)^2/\Var J\ge0$. Also
$\E(Z-\mathcal PZ)^2\le\E(Z-G)^2$ for every $G$ in the span of $1$ and $X^2$. We
take
\[
 G=y\Bigl(1-\frac{W^2}2\Bigr)+\frac12 .
\]
From now on $0<q\le\frac12$ and $1\le x\le N^{1/4}$, and $C$ denotes an
absolute constant that may change from line to line. Then $p\ge\frac12$,
$x\le y\le2x$, and $\Var J=(N-1)pq\ge x/4$ for $N\ge2$. So it is enough to
show that $\E(Z-G)^2\le C(x^{-2}+x^4/N)$.

For an interior bin $0<k<N$ put $s_k=\sqrt{k(N-k)}$, $c_k=N/(2s_k)$ and
$z_k=s_ku$. On the event $K=k$ we have $k(N-k)=N^2(1-W^2)/4$, so
$c_k=(1-W^2)^{-1/2}\ge1$, $2z_k=y\sqrt{1-W^2}$ and $c_kz_k=y/2$. Let
$\mathcal I_k(z)=I_0(2z)+c_kI_1(2z)$ and
$\widetilde S_k(v)=2v\,\mathcal I_k(s_kv)$. This is Paper~A's Bessel
approximation to $S_{k,N-k}(v)$, which is called $T_{a,b}$
in~\cite{paperA}. By~\cite[Thm.~6.2]{paperA},
\begin{equation}\label{eq:fisher-bessel-bound}
 0\le1-\frac{S_{k,N-k}(v)}{\widetilde S_k(v)}\le\frac{Nv^2}2+v
 \qquad\text{for all }v>0 .
\end{equation}
Let $\Lambda_k(t)=\log S_{k,N-k}(e^t)$ and
$\beta_k(t)=\log\widetilde S_k(e^t)$. By Proposition~\ref{prop:eta},
$Z=\Lambda_K'(\log u)$ on the event $0<K<N$.

\emph{Step 1.} We compare $\Lambda_k'$ with $\beta_k'$ at $\log u$. Both
functions are convex, since each is the logarithm of a series in $e^t$ with
nonnegative coefficients. Let $0<\delta\le1$ and $\lvert t-\log u\rvert\le
\delta$, so that $v=e^t\le eu$. Then
\[
 \frac{Nv^2}2+v\le\frac{e^2y^2/2+ey}N\le\frac{2e^2+2e}{\sqrt N},
\]
using $y\le2N^{1/4}$. The right side is at most $\frac12$ for $N\ge1700$.
Since $-\log(1-\varepsilon)\le2\varepsilon$ for
$0\le\varepsilon\le\frac12$, the bound \eqref{eq:fisher-bessel-bound} gives
$0\le\beta_k-\Lambda_k\le\eta$ on this interval, with
$\eta=(e^2y^2+2ey)/N$.

Next, $\beta_k(t)=\log2+t+\log\mathcal I_k(s_ke^t)$, and
$\mathcal I_k(s_ke^t)=\sum_na_ne^{nt}$ with $a_n\ge0$. So $\beta_k''(t)$ is
the variance of $n$ under the weights $a_ne^{nt}$, which is at most the mean
of $n^2$. From the power series,
$(z\frac d{dz})^2I_0(2z)=4z^2I_0(2z)$ and
$(z\frac d{dz})^2I_1(2z)=(4z^2+1)I_1(2z)$. So the mean of $n^2$ is at most
$4z^2+1$ with $z=s_ke^t\le ey/2$, since $s_k\le N/2$. Hence
$\beta_k''\le\kappa=e^2y^2+1$ on the interval.

Now take $\delta=\sqrt{2\eta/\kappa}$. It is at most $1$, since for $N\ge4$
we have $2\eta\le\frac12(e^2y^2+2ey)\le\kappa$. Lemma~\ref{lem:fisher-convex}
with $f=\Lambda_k$ and $g=\beta_k$ gives, for every $0<k<N$,
\begin{equation}\label{eq:fisher-step1}
 \bigl\lvert\Lambda_k'(\log u)-\beta_k'(\log u)\bigr\rvert
 \le\frac\eta\delta+\frac{\kappa\delta}2=\sqrt{2\eta\kappa}
 \le\frac{2e^2y^2+2ey+1}{\sqrt{2N}}\le\frac{11(1+y)^2}{\sqrt N} .
\end{equation}
The second inequality is $\sqrt{ab}\le\frac12(a+b)$.

\emph{Step 2.} We compute $\beta_k'(\log u)$. Write $z=s_ke^t$, $w=2z$ and
$R=R(w)$ as in Lemma~\ref{lem:fisher-bessel}. Using $I_0'=I_1$ and
$I_1'(w)=I_0(w)-I_1(w)/w$,
\[
 \beta_k'(t)=1+\frac{z\mathcal I_k'(z)}{\mathcal I_k(z)}
 =1+\frac{wR+c_kw-c_kR}{1+c_kR}=w+\frac{1+w(c_k-1)(1-R)}{1+c_kR}.
\]
At $t=\log u$ we have $w=2z_k$. Call the last fraction $B_k$ at this point.
So $\beta_k'(\log u)=2z_k+B_k$. By Lemma~\ref{lem:fisher-bessel},
$0<R<1$ and $w(1-R)\le1$. Since $c_k\ge1$, the numerator of $B_k$ lies
between $1$ and $c_k$. The numerator is also at most $1+wc_k=1+y$, and the
denominator is at least $1$, so $0<B_k\le1+y$. Moreover
\[
 B_k-\frac12\le\frac{c_k}{1+c_kR}-\frac12
 =\frac{(c_k-1)+c_k(1-R)}{2(1+c_kR)}\le\frac{c_k-1}2+\frac{c_k(1-R)}2 ,
\]
and $B_k-\frac12\ge\frac{1-c_kR}{2(1+c_kR)}\ge-\frac{c_k-1}2$. By
Lemma~\ref{lem:fisher-bessel}, $c_k(1-R)\le c_k/w=c_k^2/y$. Hence
\[
 \Bigl\lvert B_k-\frac12\Bigr\rvert\le\frac{c_k-1}2+\frac{c_k^2}{2y}.
\]
If $W^2\le\frac12$ on $K=k$, then $c_k^2\le2$ and $c_k-1\le W^2$, because
$(1-v)^{-1/2}-1$ is convex in $v$, vanishes at $v=0$ and is less than
$\frac12$ at $v=\frac12$, so it is at most $v$ for $0\le v\le\frac12$. Hence
$\lvert B_k-\frac12\rvert\le W^2/2+1/y$ there.

The term $B_k$ is the part of the score that does not grow with $x$. It is
close to $\frac12$ in the middle of the row, and this is the $\frac12$ in
$G$.

\emph{Step 3.} Let $\psi(w)=1-w^2/2-\sqrt{1-w^2}$ for $\lvert w\rvert\le1$.
The power series of $\sqrt{1-v}$ is $1-v/2$ minus a series
$\sum_{j\ge2}a_jv^j$ with $a_j\ge0$, and $\sum_ja_j=\frac12$ since
$\sqrt{1-v}$ vanishes at $v=1$. So $0\le\psi(w)\le w^4/2$. On $K=k$ we have
$2z_k=y\sqrt{1-W^2}=y(1-W^2/2)-y\psi(W)$. So on the event $0<K<N$, Steps~1
and~2 give
\[
 Z-G=\bigl(Z-\beta_K'(\log u)\bigr)+\Bigl(B_K-\frac12\Bigr)-y\psi(W).
\]
On the 2 atoms $K\in\{0,N\}$ we have $Z=0$ and $W^2=1$, so
$Z-G=-(y+1)/2$.

\emph{Step 4.} We collect the moments we need. By
Lemma~\ref{lem:fisher-W}, $\E W^4\le3/x^2+9/(2x^4)\le8/x^2$ and
$\E W^8\le Cx^{-4}$. By Markov's inequality,
$\Prob(W^2>\frac12)\le16\,\E W^8\le Cx^{-4}$. Also
$\Prob(K\in\{0,N\})=\Prob(J=0)=p^{N-1}\le e^{-q(N-1)}\le e^{1/2}e^{-x}$.

\emph{Step 5.} Write $\lVert\cdot\rVert$ for the norm of $L^2$ and
$\mathbf 1_{\rm int}$ for the indicator of $0<K<N$. By Step~3 and
Minkowski's inequality,
\[
 \lVert Z-G\rVert\le
 \bigl\lVert(Z-\beta_K'(\log u))\mathbf 1_{\rm int}\bigr\rVert
 +\bigl\lVert(B_K-\tfrac12)\mathbf 1_{\rm int}\bigr\rVert
 +\lVert y\psi(W)\rVert
 +\tfrac{y+1}2\,\Prob(K\in\{0,N\})^{1/2}.
\]
By \eqref{eq:fisher-step1} the first term is at most
$11(1+y)^2/\sqrt N\le99x^2/\sqrt N$, since $(1+y)^2\le(1+2x)^2\le9x^2$. For
the second term we use Step~2 on the event $W^2\le\frac12$, and
$\lvert B_K-\frac12\rvert\le1+y$ on the rest of the interior. This gives
\[
 \E\bigl[(B_K-\tfrac12)^2\mathbf 1_{\rm int}\bigr]
 \le\E\Bigl(\frac{W^2}2+\frac1y\Bigr)^2+(1+y)^2\Prob\Bigl(W^2>\frac12\Bigr)
 \le\frac{\E W^4}2+\frac2{y^2}+\frac{C(1+y)^2}{x^4}\le\frac C{x^2}.
\]
By Step~3 the third term is at most $\frac y2(\E W^8)^{1/2}\le C/x$. By
Step~4 the last term is at most $Cxe^{-x/2}\le C/x$. Therefore
\[
 \E(Z-G)^2\le C\Bigl(\frac{x^4}N+\frac1{x^2}\Bigr),
\]
and dividing by $\Var J\ge x/4$ gives
$e_N-R^2_N\le C(x^{-3}+x^3/N)$.

We used $N\ge1700$ only in Step~1, so we may take $N_0=1700$. Every other
constant is absolute, since each bound above holds uniformly in
$0<q\le\frac12$ and $1\le x\le N^{1/4}$. This proves the bound in
Theorem~\ref{thm:fisher}. The expansion of $e_N$ in
Theorem~\ref{thm:fisher}, Corollary~\ref{cor:fisher-crossing} and the lower
bound in Remark~\ref{rem:fisher-regimes} follow from this bound and the
closed form, and their proofs are given in Section~\ref{sec:fisher}.
